\SourceLeaf{409}%
\[
aa=\frac{4\sin\frac32\omega.\sin A^2}{\sin(A+\frac12\omega).\sin(B+\frac12\omega).\sin(C+\frac12\omega)}
\]
\[
\log\frac{a\sin B}{b\sin A}=\frac{(bb-aa)(4aa+4bb-5cc)}{720}
\]
\[
\begin{aligned}
\log\frac{3\omega}{\frac12 ab\sin C}&=\tfrac12\omega\{\operatorname{cotg}A+\operatorname{cotg}B+\operatorname{cotg}C\}\\
&=\tfrac12\omega\frac{1+\cos A\cos B\cos C}{\sin A\sin B\sin C}.
\end{aligned}
\]
\par\Needspace{9\baselineskip}\begin{center}[2.]\end{center}
\noindent\hspace*{1.2em}Die allgemeine Untersuchung der Dreiecke auf krummen Flächen hat uns folgendes wichtige Resultat gegeben.

Es seien $\BareRoot{\frac1\alpha}$, $\BareRoot{\frac1\beta}$, $\BareRoot{\frac1\gamma}$ die Halbmesser der Kugeln, die dieselbe Krümmung haben wie die krumme Fläche resp. an den Punkten $A,B,C$: dann ist
\[
\begin{aligned}
\log\frac{b\sin C}{c\sin B}&=\tfrac1{24}(bb-cc)\alpha+\tfrac1{48}(-aa+3bb-3cc)\beta\\
&\quad+\tfrac1{48}(aa+3bb-3cc)\gamma\\
A+B+C&=180^0+\tfrac16(\alpha+\beta+\gamma)bc\sin A,
\end{aligned}
\]
\noindent und wenn man
\[
\begin{aligned}
dA&=\frac{bc\sin A}{24}(2\alpha+\beta+\gamma)\\
dB&=\frac{bc\sin A}{24}(\alpha+2\beta+\gamma)\\
dC&=\frac{bc\sin A}{24}(\alpha+\beta+2\gamma)
\end{aligned}
\]
\noindent macht, so wird
\[
\frac{\sin(A-dA)}a=\frac{\sin(B-dB)}b=\frac{\sin(C-dC)}c.
\]
\indent Es ist unnöthig zu bemerken, dass hier $A,B,C$ die Winkel des sphäroidischen Dreiecks selbst bedeuten.

Hiernach finden sich z. B. in unserm grössten Dreiecke zwischen Hohe\-hagen---Brocken---Inselsberg die Winkelcorrectionen
\begin{center}\begin{tabular}{lrr}
Inselsberg & $-\overset{\prime\prime}{4}{,}95247$ & $\overset{\prime\prime}{4}{,}95119$\\
Hohehagen & $-4{,}95264$ & $4{,}95102$\\
Brocken & $-4{,}95273$ & $4{,}95093$\\
Summe & $-14{,}85784$ & $14{,}85314.$
\end{tabular}\end{center}
